\documentclass[11pt,a4paper]{article}
\usepackage{fontspec}
\usepackage[ngerman]{babel}
\setmainfont{Latin Modern Roman}
\setsansfont{Latin Modern Sans}
\usepackage{amsmath,amssymb,mathtools}
\usepackage[a4paper,left=2cm,right=4.3cm,top=2.3cm,bottom=2.4cm,headheight=14pt]{geometry}
\usepackage{tikz}
\usetikzlibrary{arrows.meta,calc,decorations.markings}
\usepackage[most]{tcolorbox}
\usepackage{booktabs,array,tabularx,enumitem,siunitx,titlesec,fancyhdr,caption}
\usepackage[hidelinks]{hyperref}
\sisetup{locale=DE}

\definecolor{acc}{RGB}{31,95,170}
\definecolor{acc2}{RGB}{205,85,40}
\definecolor{acc3}{RGB}{30,140,90}
\definecolor{merk}{RGB}{225,145,25}
\definecolor{wasser}{RGB}{214,232,250}

\titleformat{\section}{\sffamily\Large\bfseries\color{acc}}{\thesection}{0.7em}{}
\titleformat{\subsection}{\sffamily\large\bfseries}{\thesubsection}{0.6em}{}
\titleformat{\subsubsection}{\sffamily\normalsize\bfseries}{}{0em}{}
\titlespacing*{\section}{0pt}{1.2em}{0.6em}
\titlespacing*{\subsection}{0pt}{1em}{0.4em}

\pagestyle{fancy}
\fancyhf{}
\fancyhead[L]{\sffamily\small Strömungslehre -- Lösungen der Übungsaufgaben}
\fancyhead[R]{\sffamily\small\thepage}
\renewcommand{\headrulewidth}{0.3pt}

\captionsetup{font={small,sf},labelfont=bf,skip=4pt}
\setlist{itemsep=1pt,topsep=3pt}
\setlength{\parindent}{0pt}
\setlength{\parskip}{0.45em}
\setlength{\emergencystretch}{2em}

\newtcolorbox{formel}[1][]{colback=acc!5,colframe=acc!75!black,boxrule=0.6pt,arc=2pt,
  left=5pt,right=5pt,top=3pt,bottom=3pt,fonttitle=\sffamily\small\bfseries,#1}
\newtcolorbox{merke}[1][Merke]{colback=merk!9,colframe=merk!85!black,boxrule=0.6pt,arc=2pt,
  left=5pt,right=5pt,top=3pt,bottom=3pt,title={#1},fonttitle=\sffamily\small\bfseries}

\tikzset{>=Stealth,
  wand/.style={line width=1.2pt},
  fluid/.style={fill=wasser},
  pf/.style={->,line width=1.1pt},
  mass/.style={<->,thin},
  lbl/.style={font=\small}}

\newcommand{\abb}[1]{\captionof{figure}{#1}}
\newenvironment{bild}{\par\medskip\noindent\begin{minipage}{\linewidth}\centering}{\end{minipage}\par\medskip}
\newcommand{\dd}{\mathrm{d}}


\newtcolorbox{ergebnis}[1][Ergebnis]{colback=acc3!7,colframe=acc3!80!black,boxrule=0.6pt,arc=2pt,
  left=5pt,right=5pt,top=3pt,bottom=3pt,title={#1},fonttitle=\sffamily\small\bfseries}
\newcommand{\gegeben}{\textbf{Gegeben:}\ }
\newcommand{\gesucht}{\textbf{Gesucht:}\ }

\begin{document}

\begin{center}
{\sffamily\bfseries\huge\color{acc} Grundlagen der Strömungslehre}\\[4pt]
{\sffamily\Large Lösungen der Übungsaufgaben}\\[6pt]
{\small Ausführliche Lösungswege zum Aufgabenblatt der Vorlesung von Prof.\ Dr.-Ing.\ G.\ Thieleke}
\end{center}

\vspace{4pt}
Jede Lösung folgt demselben Schema: Skizze, Ansatz (welches Werkzeug und warum), Rechnung mit Zwischenergebnissen, Ergebnis und eine kurze Kontrolle. Gerechnet ist mit $g = \SI{9.81}{m/s^2}$ und, wo nicht anders angegeben, $\rho_{Wasser} = \SI{1000}{kg/m^3}$. Die Ergebnisse sind nachgerechnet, aber nicht mit offiziellen Musterlösungen abgeglichen.

\begin{merke}[Umfang]
Dieses Heft enthält die Aufgaben 1, 2, 7, 8, 9 und 13. Für die übrigen Aufgaben des Blatts lag der Aufgabentext nicht vor; sie werden ergänzt, sobald das Aufgabenblatt vorliegt.
\end{merke}

\tableofcontents

%=====================================================================
\section{Aufgabe 1 -- Druckluft-Düse}

\begin{bild}
% Aufgabe 1: Kessel mit Duese
\begin{tikzpicture}[scale=0.85]
\fill[wasser] (0,0) rectangle (3,3);
\fill[wasser] (3,1.0) .. controls (4,1.0) and (4.4,1.3) .. (5.6,1.3) -- (5.6,1.7) .. controls (4.4,1.7) and (4,2.0) .. (3,2.0) -- cycle;
\draw[wand] (3,2.0) -- (3,3) -- (0,3) -- (0,0) -- (3,0) -- (3,1.0);
\draw[wand] (3,2.0) .. controls (4,2.0) and (4.4,1.7) .. (5.6,1.7);
\draw[wand] (3,1.0) .. controls (4,1.0) and (4.4,1.3) .. (5.6,1.3);
\node[lbl,align=center] at (1.5,1.5) {Kessel\\$p_0,\ T_0$\\$v_0 \approx 0$};
\draw[thin,dashed] (5.6,0.9)--(5.6,2.1) node[above]{1};
\draw[pf,acc2] (5.8,1.5)--(7.0,1.5) node[right,black]{$v_1$};
\node[lbl,anchor=north] at (5.6,0.85) {$p_1,\ T_1,\ \rho_1$};
\end{tikzpicture}

\abb{Druckluft strömt aus einem Kessel durch eine Düse ins Freie.}
\end{bild}

\gegeben Kessel $p_0 = \SI{6}{bar}$, $T_0 = \SI{313.15}{K}$ (\SI{40}{\celsius}), $v_0 \approx 0$; Austrittsdruck $p_1 = \SI{4}{bar}$; Luft mit $R = \SI{287.1}{J/(kg.K)}$, $\kappa = 1{,}4$.

\gesucht a) $T_1$, $v_1$, $\rho_1$ bei verlustfreier Strömung; b) $v_1$, wenn man die Luft als inkompressibel rechnet, und der Fehler; c) Zustand bei der gemessenen Geschwindigkeit $v_{1c} = \SI{220}{m/s}$ und der Düsenwirkungsgrad.

\textbf{Ansatz:} Bei Gasen mit hoher Geschwindigkeit ändert sich die Dichte merklich. Deshalb wird nicht mit Bernoulli, sondern mit dem Energiesatz für Gase gerechnet:
\[
\frac{v_0^2}{2} + c_p T_0 = \frac{v_1^2}{2} + c_p T_1,
\qquad
c_p = \frac{\kappa}{\kappa - 1}\,R = \frac{1{,}4}{0{,}4}\cdot 287{,}1 = \SI{1004.9}{J/(kg.K)}
\]

\subsection*{a) Verlustfrei (isentrop)}

Isentrope zwischen Kessel und Austritt, Exponent $\frac{\kappa - 1}{\kappa} = 0{,}2857$:
\[
T_1 = T_0\left(\frac{p_1}{p_0}\right)^{\frac{\kappa-1}{\kappa}} = 313{,}15\cdot\left(\frac{4}{6}\right)^{0{,}2857}\,\mathrm{K} = 313{,}15\cdot 0{,}8906\,\mathrm{K} = \SI{278.9}{K}\ (\SI{5.7}{\celsius})
\]
Energiesatz mit $v_0 = 0$:
\[
v_1 = \sqrt{2\,c_p\,(T_0 - T_1)} = \sqrt{2\cdot 1004{,}9\cdot 34{,}25}\ \mathrm{m/s} = \sqrt{68\,840}\ \mathrm{m/s} = \SI{262.4}{m/s}
\]
Gasgleichung:
\[
\rho_1 = \frac{p_1}{R\,T_1} = \frac{4\cdot 10^5}{287{,}1\cdot 278{,}9}\ \mathrm{kg/m^3} = \SI{5.00}{kg/m^3}
\]
\textbf{Kontrolle:} $p_1/p_0 = 0{,}667 > 0{,}528$, das Druckverhältnis ist unterkritisch. Schallgeschwindigkeit $a_1 = \sqrt{\kappa R T_1} = \SI{335}{m/s}$, also $Ma_1 = 0{,}78 < 1$. Eine einfache (konvergente) Düse reicht.

\subsection*{b) Inkompressibel gerechnet}

Dichte bleibt auf dem Kesselwert:
\[
\rho_0 = \frac{p_0}{R\,T_0} = \frac{6\cdot 10^5}{287{,}1\cdot 313{,}15}\ \mathrm{kg/m^3} = \SI{6.67}{kg/m^3}
\]
Bernoulli zwischen Kessel und Austritt:
\[
v_1 = \sqrt{\frac{2\,(p_0 - p_1)}{\rho_0}} = \sqrt{\frac{2\cdot 2\cdot 10^5}{6{,}674}}\ \mathrm{m/s} = \SI{244.8}{m/s}
\]
\[
\text{Fehler} = \frac{244{,}8 - 262{,}4}{262{,}4} = \SI{-6.7}{\percent}
\]
Die Faustregel ($\rho \approx \text{const}$ nur bis etwa \SI{70}{m/s}) ist hier klar verletzt.

\subsection*{c) Mit Verlusten, $v_{1c} = \SI{220}{m/s}$}

Die Düse ist adiabat, der Energiesatz gilt also weiter. Nur die Isentrope gilt nicht mehr.
\[
T_{1c} = T_0 - \frac{v_{1c}^2}{2\,c_p} = 313{,}15\,\mathrm{K} - \frac{220^2}{2\cdot 1004{,}9}\,\mathrm{K} = (313{,}15 - 24{,}08)\,\mathrm{K} = \SI{289.1}{K}\ (\SI{15.9}{\celsius})
\]
\[
\rho_{1c} = \frac{p_1}{R\,T_{1c}} = \frac{4\cdot 10^5}{287{,}1\cdot 289{,}1}\ \mathrm{kg/m^3} = \SI{4.82}{kg/m^3}
\]
Düsenwirkungsgrad als Verhältnis der kinetischen Energien (wirklich zu isentrop):
\[
\eta_D = \frac{v_{1c}^2/2}{v_1^2/2} = \left(\frac{220}{262{,}4}\right)^2 = 0{,}70
\]
Die Reibung erwärmt das Gas: $T_{1c} > T_1$ bei gleichem Druck, deshalb ist auch die Dichte kleiner.

\begin{ergebnis}
\begin{tabular}{@{}ll@{}}
a) & $T_1 = \SI{278.9}{K}$, \ $v_1 = \SI{262}{m/s}$, \ $\rho_1 = \SI{5.0}{kg/m^3}$ \\
b) & $v_1 = \SI{245}{m/s}$, Fehler $\approx \SI{-7}{\percent}$ \\
c) & $T_{1c} = \SI{289.1}{K}$, \ $\rho_{1c} = \SI{4.82}{kg/m^3}$, \ $\eta_D = 0{,}70$
\end{tabular}
\end{ergebnis}

%=====================================================================
\section{Aufgabe 2 -- Gashähne im Hochhaus}

\begin{bild}
% Aufgabe 2: Gasleitung im Hochhaus
\begin{tikzpicture}[scale=0.85]
\draw[thick,gray] (0,0) rectangle (3.2,5);
\foreach \y in {1,2,3,4} \draw[thin,gray!60] (0,\y)--(3.2,\y);
\fill[wasser] (1.45,-0.6) rectangle (1.75,4.5);
\draw[wand] (1.45,-0.6)--(1.45,4.5)--(1.75,4.5)--(1.75,-0.6);
\draw[thick] (1.75,0.4)--(2.4,0.4); \fill (2.4,0.4) circle (1.6pt);
\draw[thick] (1.75,4.2)--(2.4,4.2); \fill (2.4,4.2) circle (1.6pt);
\draw[pf,acc2] (2.5,0.4)--(3.6,0.4) node[right,black]{$v_{unten}$};
\draw[pf,acc2] (2.5,4.2)--(3.6,4.2) node[right,black]{$v_{oben}$};
\node[lbl,anchor=west] at (3.65,0.0) {$\Delta p_{unten} = \SI{2000}{Pa}$};
\node[lbl,anchor=west] at (3.65,3.8) {$\Delta p_{oben} = ?$};
\draw[mass] (-0.4,0.4)--(-0.4,4.2) node[midway,left]{$h = \SI{30}{m}$};
\node[lbl] at (1.6,-0.9) {Gas $\rho_G$};
\node[lbl,anchor=west] at (3.65,2.1) {Luft $\rho_L$};
\end{tikzpicture}

\abb{Steigleitung mit Gas; ein Hahn im Erdgeschoss, einer \SI{30}{m} höher.}
\end{bild}

\gegeben Überdruck am unteren Hahn $\Delta p_{unten} = \SI{2000}{Pa}$, Höhenunterschied $h = \SI{30}{m}$, Dichtedifferenz Luft -- Gas $\rho_L - \rho_G = \SI{0.7}{kg/m^3}$. Das Gas ruht in der Leitung.

\gesucht Verhältnis der Ausströmgeschwindigkeiten $v_{oben}/v_{unten}$.

\textbf{Ansatz:} Hydrostatik getrennt für die ruhende Gassäule in der Leitung und die Luft außen. Über \SI{30}{m} sind die Dichten konstant (Höhenformel nicht nötig).
\[
p_{G,oben} = p_{G,unten} - \rho_G\,g\,h,
\qquad
p_{L,oben} = p_{L,unten} - \rho_L\,g\,h
\]
Subtrahieren liefert den Überdruck oben:
\[
\Delta p_{oben} = p_{G,oben} - p_{L,oben} = \Delta p_{unten} + (\rho_L - \rho_G)\,g\,h = 2000\,\mathrm{Pa} + 0{,}7\cdot 9{,}81\cdot 30\,\mathrm{Pa} = \SI{2206}{Pa}
\]
Gas ist leichter als Luft, sein Druck nimmt nach oben langsamer ab -- der Überdruck wächst mit der Höhe.

Ausströmen aus dem Hahn (Bernoulli, Gas in der Leitung ruht): $v = \sqrt{2\,\Delta p/\rho_G}$. Die Gasdichte kürzt sich heraus:
\[
\frac{v_{oben}}{v_{unten}} = \sqrt{\frac{\Delta p_{oben}}{\Delta p_{unten}}} = \sqrt{\frac{2206}{2000}} = 1{,}05
\]

\begin{ergebnis}
$\Delta p_{oben} \approx \SI{2206}{Pa}$, \quad $v_{oben}/v_{unten} \approx 1{,}05$ -- oben strömt das Gas um rund \SI{5}{\percent} schneller aus.
\end{ergebnis}

%=====================================================================
\section{Aufgabe 7 -- Körper im Rohr}

\begin{bild}
% Aufgabe 7: Koerper im Rohr
\begin{tikzpicture}[scale=0.85]
\fill[wasser] (0,-1.5) rectangle (10,1.5);
\draw[wand] (0,1.5)--(10,1.5); \draw[wand] (0,-1.5)--(10,-1.5);
\fill[gray!55] (3.0,-0.75) .. controls (2.2,-0.75) and (2.2,0.75) .. (3.0,0.75) -- (3.8,0.75) -- (3.8,-0.75) -- cycle;
\draw[thick] (3.0,-0.75) .. controls (2.2,-0.75) and (2.2,0.75) .. (3.0,0.75) -- (3.8,0.75) -- (3.8,-0.75) -- cycle;
\node[lbl] at (3.25,0) {$A_K$};
\draw[acc,thick] (3.8,0.75) .. controls (5.6,0.75) and (6.6,0.3) .. (7.6,0.0);
\draw[acc,thick] (3.8,-0.75) .. controls (5.6,-0.75) and (6.6,-0.3) .. (7.6,0.0);
\foreach \x in {0.6,1.0} \draw[thin,gray] (4.6,0) ellipse ({\x} and {0.45*\x});
\draw[thin,dashed] (1.0,-1.8)--(1.0,1.8) node[above]{1};
\draw[thin,dashed] (3.8,-1.8)--(3.8,1.8) node[above]{2};
\draw[thin,dashed] (9.0,-1.8)--(9.0,1.8) node[above]{3};
\foreach \y in {-1.0,0,1.0} \draw[->,thick,acc2] (0.15,\y)--(0.9,\y);
\foreach \y in {1.12,-1.12} \draw[->,thick,acc2] (3.9,\y)--(5.2,\y);
\foreach \y in {-1.0,0,1.0} \draw[->,thick,acc2] (9.1,\y)--(9.85,\y);
\node[lbl,fill=wasser,inner sep=1pt] at (1.9,-0.6) {$v_1, p_1$};
\node[lbl,fill=wasser,inner sep=1pt] at (6.0,1.12) {$v_2, p_2$};
\node[lbl,fill=wasser,inner sep=1pt] at (8.25,-0.6) {$v_3, p_3$};
\node[lbl] at (8.3,1.15) {$A$};
\end{tikzpicture}

\abb{Ein Körper mit Stirnfläche $A_K$ verengt das Rohr; dahinter vermischt sich der Strahl wie bei der plötzlichen Erweiterung.}
\end{bild}

\gegeben Wasser, $v_1 = \SI{10}{m/s}$, $p_1 = \SI{2}{bar}$, Flächenverhältnis $\alpha = A_K/A = 0{,}5$. Keine Wandreibung.

\gesucht $p_2$ am Körperende, $p_3$ nach der Vermischung, Widerstandsbeiwert $c_W$ des Körpers.

\textbf{Ansatz:} $1 \to 2$: Beschleunigung im Spalt, verlustfrei $\Rightarrow$ Kontinuität und Bernoulli. $2 \to 3$: Vermischung hinter dem Körper $\Rightarrow$ Impulssatz (wie Carnot). $1 \to 3$: Kraft auf den Körper $\Rightarrow$ Impulssatz.

\subsection*{Schritt 1: Querschnitt 2}
\[
v_2 = v_1\,\frac{A}{A - A_K} = \frac{v_1}{1 - \alpha} = \frac{10}{0{,}5}\,\mathrm{m/s} = \SI{20}{m/s}
\]
\[
p_2 = p_1 + \frac{\rho}{2}\left(v_1^2 - v_2^2\right) = 2\cdot 10^5\,\mathrm{Pa} + 500\cdot(100 - 400)\,\mathrm{Pa} = \SI{0.5e5}{Pa} = \SI{0.5}{bar}
\]

\subsection*{Schritt 2: Querschnitt 3 (Impulssatz $2 \to 3$)}

Wie bei Carnot wirkt $p_2$ auch auf die Rückseite des Körpers, also auf den ganzen Querschnitt $A$. Durch den Spalt tritt der Impulsstrom $\rho\,v_2^2\,(1 - \alpha)A$ ein, in 3 ist $v_3 = v_1$ (Kontinuität):
\[
p_2\,A + \rho\,v_2^2\,(1 - \alpha)\,A = p_3\,A + \rho\,v_3^2\,A
\]
\[
p_3 = p_2 + \rho\left[v_2^2(1 - \alpha) - v_3^2\right] = 0{,}5\cdot 10^5\,\mathrm{Pa} + 1000\cdot(400\cdot 0{,}5 - 100)\,\mathrm{Pa} = \SI{1.5}{bar}
\]

\subsection*{Schritt 3: Widerstandskraft und $c_W$}

Impulssatz von 1 bis 3: Die Impulsströme heben sich auf ($v_1 = v_3$), es bleibt
\[
F_W = (p_1 - p_3)\,A = \SI{0.5}{bar}\cdot A
\]
Bezogen auf Staudruck und Stirnfläche des Körpers:
\[
c_W = \frac{F_W}{\frac{\rho}{2}\,v_1^2\,A_K} = \frac{p_1 - p_3}{\frac{\rho}{2}\,v_1^2\,\alpha} = \frac{0{,}5\cdot 10^5}{500\cdot 100\cdot 0{,}5} = 2{,}0
\]
Allgemein mit $p_1 - p_3 = \frac{\rho}{2}(v_2 - v_1)^2 = \frac{\rho}{2}v_1^2\left(\frac{\alpha}{1-\alpha}\right)^2$:
\[
c_W = \frac{\alpha}{(1 - \alpha)^2}
\]

\textbf{Kontrolle:} Der Druckverlust $p_1 - p_3 = \SI{0.5}{bar}$ ist genau der Carnot-Verlust $\frac{\rho}{2}(v_2 - v_3)^2 = 500\cdot 10^2\,\mathrm{Pa}$. \checkmark

\begin{ergebnis}
$p_2 = \SI{0.5}{bar}$, \quad $p_3 = \SI{1.5}{bar}$, \quad $c_W = \dfrac{\alpha}{(1-\alpha)^2} = 2{,}0$
\end{ergebnis}

%=====================================================================
\section{Aufgabe 8 -- Sieb im Rohr}

\begin{bild}
% Aufgabe 8: Sieb im Rohr
\begin{tikzpicture}[scale=0.85]
\fill[wasser] (0,-1.4) rectangle (8,1.4);
\draw[wand] (0,1.4)--(8,1.4); \draw[wand] (0,-1.4)--(8,-1.4);
\draw[thick,dash pattern=on 2pt off 1.5pt] (4,-1.4)--(4,1.4);
\draw[thick] (3.92,-1.4)--(3.92,1.4); \draw[thick] (4.08,-1.4)--(4.08,1.4);
\node[lbl,anchor=south] at (4,1.45) {Sieb};
\draw[thin,dashed] (1.0,-1.7)--(1.0,1.7) node[above]{1};
\draw[thin,dashed] (6.8,-1.7)--(6.8,1.7) node[above]{2};
\draw[thin] (1.0,-1.4)--(1.0,1.4);
\draw[thick,acc2,domain=-1.4:1.4,samples=30,smooth] plot ({1.0+2.0*(1-(\x/1.4)^2)},\x);
\foreach \y in {-1.0,-0.5,0,0.5,1.0} \draw[->,thin,acc2] (1.0,\y)--({1.0+2.0*(1-(\y/1.4)^2)},\y);
\node[lbl,fill=wasser,inner sep=1pt] at (2.6,-1.1) {$v(r)$};
\foreach \y in {-1.0,-0.5,0,0.5,1.0} \draw[->,thin,acc2] (6.8,\y)--(7.8,\y);
\node[lbl,anchor=west] at (8.4,0.5) {$\bar v$};
\node[lbl,anchor=north] at (1.0,-1.75) {$p_1$};
\node[lbl,anchor=north] at (6.8,-1.75) {$p_2$};
\draw[mass] (-0.3,0)--(-0.3,1.4) node[midway,left]{$R$};
\draw[thin,dash dot] (-0.5,0)--(8.3,0);
\draw[pf,acc] (4,-1.5)--(4,-2.3) node[right,black]{$\dot Q$};
\end{tikzpicture}

\abb{Vor dem Sieb laminares Parabelprofil, dahinter gleichmäßiges Profil. Die im Sieb verlorene Energie wird als Wärme $\dot Q$ abgeführt.}
\end{bild}

\gegeben Rohr mit Radius $R$, Drücke $p_1$, $p_2$, Dichte $\rho$; in 1 laminares Profil $v(r) = v_{max}\left(1 - r^2/R^2\right)$, in 2 gleichmäßiges Profil.

\gesucht a) Kraft $F$ auf das Sieb, b) abzuführender Wärmestrom $\dot Q$ (Temperatur soll konstant bleiben).

\textbf{Ansatz:} Weil das Profil in 1 nicht gleichmäßig ist, braucht man die Korrekturbeiwerte für Impuls ($\beta$) und kinetische Energie ($\alpha$).

\subsection*{Vorarbeit: mittlere Geschwindigkeit und Beiwerte}
\[
\bar v = \frac{1}{\pi R^2}\int_0^R v_{max}\left(1 - \frac{r^2}{R^2}\right) 2\pi r\,\dd r = \frac{v_{max}}{2} = v_2
\]
\[
\beta = \frac{1}{A}\int_A\left(\frac{v}{\bar v}\right)^2 \dd A = \frac{v_{max}^2\,\pi R^2/3}{\bar v^2\,\pi R^2} = \frac43,
\qquad
\alpha = \frac{1}{A}\int_A\left(\frac{v}{\bar v}\right)^3 \dd A = \frac{v_{max}^3\,\pi R^2/4}{\bar v^3\,\pi R^2} = 2
\]

\subsection*{a) Kraft auf das Sieb (Impulssatz)}

Kontrollfläche von 1 bis 2, Wandreibung vernachlässigt:
\[
p_1 A + \beta\,\rho\,\bar v^2 A - p_2 A - \rho\,\bar v^2 A - F = 0
\]
\[
F = (p_1 - p_2)\,\pi R^2 + (\beta - 1)\,\rho\,\bar v^2\,\pi R^2 = \pi R^2\left[(p_1 - p_2) + \frac{\rho\,v_{max}^2}{12}\right]
\]
Das Sieb muss also nicht nur die Druckdifferenz aufnehmen, sondern auch den Impuls, den es beim Vergleichmäßigen des Profils aus der Strömung nimmt.

\subsection*{b) Wärmestrom (Energiebilanz)}

Bei konstanter Temperatur bleibt die innere Energie gleich. Alles, was an Druck- und kinetischer Energie verloren geht, muss als Wärme hinaus:
\[
\dot Q = Q\left[(p_1 - p_2) + (\alpha - 1)\,\rho\,\frac{\bar v^2}{2}\right],
\qquad Q = \bar v\,\pi R^2 = \frac{v_{max}}{2}\,\pi R^2
\]
\[
\dot Q = \pi R^2\,\frac{v_{max}}{2}\left[(p_1 - p_2) + \frac{\rho\,v_{max}^2}{8}\right]
\]

\begin{ergebnis}
$F = \pi R^2\left[(p_1 - p_2) + \dfrac{\rho\,v_{max}^2}{12}\right]$, \qquad
$\dot Q = \pi R^2\,\dfrac{v_{max}}{2}\left[(p_1 - p_2) + \dfrac{\rho\,v_{max}^2}{8}\right]$
\par\smallskip
Die Aufgabe nennt keine Zahlenwerte für die Geschwindigkeit, deshalb bleibt $v_{max}$ im Ergebnis.
\end{ergebnis}

%=====================================================================
\section{Aufgabe 9 -- Wasserstrahl gegen Wand im Gerinne}

\begin{bild}
% Aufgabe 9: Strahl gegen Wand im Gerinne
\begin{tikzpicture}[scale=0.85]
\fill[wasser] (0,0) -- (0,2.5) .. controls (3,2.5) and (5,2.6) .. (6,2.7) .. controls (7,2.85) and (7.5,3.6) .. (8,4.1) -- (8,0) -- cycle;
\draw[wand] (-0.2,0)--(8,0)--(8,4.6);
\draw[thick,acc] (0,2.5) .. controls (3,2.5) and (5,2.6) .. (6,2.7) .. controls (7,2.85) and (7.5,3.6) .. (8,4.1);
\fill[acc!25] (0,0) rectangle (6.3,0.35);
\draw[thin,acc] (0,0.35)--(6.3,0.35);
\draw[pf,acc2] (0.6,0.17)--(2.2,0.17); \node[font=\small,anchor=south] at (2.6,0.38) {$v_1$};
\draw[pf,acc2] (3.6,1.6)--(1.9,1.6) node[left,black,font=\small]{$v_r$};
\draw[thin,dashed] (0,-0.4)--(0,3.0);
\draw[mass] (-0.5,0)--(-0.5,2.5) node[midway,left]{$h_2$};
\draw[mass] (6.6,0)--(6.6,0.35); \node[lbl,anchor=west] at (6.65,0.25) {$h_1$};
\draw[mass] (8.5,0)--(8.5,4.1) node[midway,right]{$h$};
\node[lbl,acc,anchor=south] at (3.5,2.9) {Kontrollvolumen};
\draw[acc,densely dotted,thick] (0.05,0.05) rectangle (7.95,2.85);
\end{tikzpicture}

\abb{Ein flacher Strahl ($h_1$, $v_1$) läuft am Boden gegen eine Wand und strömt oben mit $v_r$ zurück. An der Wand staut sich das Wasser auf die Höhe $h$.}
\end{bild}

\gegeben $h_1 = \SI{0.03}{m}$, $v_1 = \SI{8}{m/s}$, Wassertiefe links $h_2 = \SI{0.5}{m}$. Breite $b$ (kürzt sich heraus).

\gesucht Wasserhöhe $h$ an der Wand.

\textbf{Ansatz:} Kontrollvolumen vom linken Schnitt (Tiefe $h_2$) bis zur Wand. In beiden Schnitten ist die Druckverteilung hydrostatisch. An der Wand ruht das Wasser.

\subsection*{Schritt 1: Rückströmung (Kontinuität)}

Was unten zuströmt, muss oben zurück:
\[
v_1\,h_1 = v_r\,(h_2 - h_1)
\quad\Rightarrow\quad
v_r = \frac{v_1\,h_1}{h_2 - h_1} = \frac{8\cdot 0{,}03}{0{,}47}\,\mathrm{m/s} = \SI{0.51}{m/s}
\]

\subsection*{Schritt 2: Impulssatz in Strömungsrichtung (pro Breite)}

Hydrostatische Druckkräfte links und an der Wand; der zuströmende Strahl und die abströmende Rückströmung bringen beide Impuls in $+x$-Richtung:
\[
\frac{\rho g}{2}\,h_2^2 + \rho\,v_1^2\,h_1 + \rho\,v_r^2\,(h_2 - h_1) = \frac{\rho g}{2}\,h^2
\]
\[
h = \sqrt{h_2^2 + \frac{2}{g}\left[v_1^2\,h_1 + v_r^2\,(h_2 - h_1)\right]}
\]

\subsection*{Schritt 3: Zahlen}
\[
v_1^2\,h_1 + v_r^2\,(h_2 - h_1) = 64\cdot 0{,}03 + 0{,}51^2\cdot 0{,}47 = (1{,}920 + 0{,}122)\,\mathrm{m^3/s^2} = \SI{2.04}{m^3/s^2}
\]
\[
h = \sqrt{0{,}25 + \frac{2\cdot 2{,}04}{9{,}81}}\ \mathrm{m} = \sqrt{0{,}666}\ \mathrm{m} = \SI{0.82}{m}
\]

\begin{ergebnis}
$v_r \approx \SI{0.51}{m/s}$, \quad $h \approx \SI{0.82}{m}$
\end{ergebnis}

%=====================================================================
\section{Aufgabe 13 -- Kreislaufprüfstand}

\begin{bild}
% Aufgabe 13: Kreislaufpruefstand (schematisch)
\begin{tikzpicture}[scale=0.85]
\draw[line width=3pt,wasser!70!acc] (0,0) rectangle (6,3.6);
% Pumpe unten Mitte, foerdert nach rechts
\draw[thick,fill=white] (3,0) circle (0.5); \draw[thick] (2.78,0.44)--(3.5,0)--(2.78,-0.44);
\node[lbl] at (3,-0.85) {Pumpe};
\fill (3.75,0) circle (1.6pt); \node[lbl,anchor=north] at (3.8,-0.08) {D};
\fill (2.25,0) circle (1.6pt); \node[lbl,anchor=north] at (2.2,-0.08) {S};
% Armatur oben Mitte
\fill[white] (2.65,3.25) rectangle (3.35,3.95);
\draw[thick] (2.65,3.25)--(3.35,3.95)--(2.65,3.95)--(3.35,3.25)--cycle;
\fill (3.7,3.6) circle (1.6pt); \node[lbl,anchor=south] at (3.75,3.68) {A};
\node[lbl,anchor=south] at (2.6,3.95) {Armatur $\zeta_A$};
% Kuehler links
\draw[thick,fill=white] (-0.35,1.45) rectangle (0.35,2.15);
\draw[thin] (-0.35,1.45)--(-0.1,1.8)--(0.1,1.5)--(0.35,1.9);
\node[lbl,anchor=west] at (0.45,1.8) {Kühler $\dot Q$};
% Kruemmer
\foreach \p in {(0,0),(6,0),(6,3.6),(0,3.6)} \fill[acc] \p circle (2.2pt);
\node[lbl,acc,anchor=west] at (6.1,3.75) {Krümmer $\zeta_{kr}$};
% Stroemungsrichtung und Leitungen
\draw[pf,acc2] (6.4,1.2)--(6.4,2.4) node[midway,right,black,font=\small]{$L_1$};
\draw[pf,acc2] (-0.7,2.9)--(-0.7,0.7) node[midway,left,black,font=\small]{$L_2$};
\draw[mass] (7.6,0)--(7.6,3.6) node[midway,right]{$h$};
\draw[thin,dotted] (6,3.6)--(7.7,3.6); \draw[thin,dotted] (6,0)--(7.7,0);
\end{tikzpicture}

\abb{Geschlossener Wasserkreislauf (schematisch): Die Pumpe fördert von D über $L_1$ mit zwei Krümmern hinauf zur Armatur (Messstelle A, Höhe $h$), von dort über $L_2$ mit zwei Krümmern und den Kühler zurück nach S.}
\end{bild}

\gegeben Wasser, $Q = \SI{0.25}{m^3/s}$, $D = \SI{0.2}{m}$, Rohrreibung $\lambda\,\frac{L_1 + L_2}{D} = 3{,}5$, Armatur $\zeta_A = 3{,}0$, je Krümmer $\zeta_{kr} = 0{,}4$, Wirkungsgrade $\eta_i = 0{,}85$, $\eta_m = 0{,}95$; Überdruck $p_A$ an der Armatur, die um $h$ über der Pumpe liegt.

\gesucht a) spezifische Förderarbeit $Y^*$, b) Antriebsleistung $P$, c) Wärmestrom $\dot Q$ im Kühler, d) Drücke $p_D$ und $p_S$ an der Pumpe.

\textbf{Annahmen:} vier Krümmer wie in der Skizze, je zwei in $L_1$ und $L_2$; $L_1$ führt von der Druckseite D zur Armatur.

\textbf{Ansatz:} Im geschlossenen Kreislauf sind Anfangs- und Endpunkt gleich. Druck, Geschwindigkeit und Höhe heben sich auf, die Pumpe deckt genau die Verluste; die Höhe $h$ fällt heraus.

\subsection*{Vorarbeit}
\[
v = \frac{Q}{\pi D^2/4} = \frac{0{,}25}{\pi\cdot 0{,}01}\,\mathrm{m/s} = \SI{7.96}{m/s},
\qquad
\frac{v^2}{2} = \SI{31.7}{J/kg},
\qquad
\rho\,\frac{v^2}{2} = \SI{0.317}{bar}
\]

\subsection*{a) Spezifische Förderarbeit}
\[
Y^* = \frac{v^2}{2}\left[\lambda\,\frac{L_1 + L_2}{D} + \zeta_A + 4\,\zeta_{kr}\right] = 31{,}7\cdot(3{,}5 + 3{,}0 + 1{,}6)\,\mathrm{J/kg} = \SI{256}{J/kg}
\]

\subsection*{b) Antriebsleistung}

Hydraulische Leistung $\rho\,Q\,Y^* = 1000\cdot 0{,}25\cdot 256\,\mathrm{W} = \SI{64.1}{kW}$. Bei der Pumpe werden die Wirkungsgrade dividiert:
\[
P = \frac{\rho\,Q\,Y^*}{\eta_i\,\eta_m} = \frac{64{,}1\,\mathrm{kW}}{0{,}85\cdot 0{,}95} = \SI{79}{kW}
\]

\subsection*{c) Wärmestrom im Kühler}

Ins Wasser geht alles außer den mechanischen Verlusten (Lager, Dichtungen): die hydraulische Leistung plus die inneren Pumpenverluste. Im stationären Betrieb muss genau das im Kühler wieder heraus:
\[
\dot Q = \eta_m\,P = \frac{\rho\,Q\,Y^*}{\eta_i} = \frac{64{,}1\,\mathrm{kW}}{0{,}85} = \SI{75}{kW}
\]

\subsection*{d) Drücke an der Pumpe}

Erweiterte Bernoulli-Gleichung mit Überdrücken, gleiche Geschwindigkeit überall. Von D hinauf zur Armatur (Höhe $h$, Leitung $L_1$, zwei Krümmer):
\[
p_D = p_A + \rho\,g\,h + \rho\,\frac{v^2}{2}\left(\lambda\,\frac{L_1}{D} + 2\,\zeta_{kr}\right) \approx \SI{4.48}{bar}
\]
Von der Armatur hinunter zur Saugseite S (Armatur, Leitung $L_2$, zwei Krümmer):
\[
p_S = p_A + \rho\,g\,h - \rho\,\frac{v^2}{2}\left(\zeta_A + \lambda\,\frac{L_2}{D} + 2\,\zeta_{kr}\right) \approx \SI{1.91}{bar}
\]
\textbf{Kontrolle:} Die Pumpe muss genau die Druckdifferenz aufbringen, die der Kreislauf verbraucht:
\[
p_D - p_S = \rho\,Y^* = 1000\cdot 256\,\mathrm{Pa} = \SI{2.56}{bar} \ \checkmark
\]

\begin{ergebnis}
\begin{tabular}{@{}ll@{}}
a) & $Y^* \approx \SI{256}{J/kg}$ \\
b) & $P \approx \SI{79}{kW}$ \\
c) & $\dot Q \approx \SI{75}{kW}$ \\
d) & $p_D \approx \SI{4.48}{bar}$, \ $p_S \approx \SI{1.91}{bar}$ (Überdruck); \ $p_D - p_S = \rho\,Y^*$
\end{tabular}
\end{ergebnis}

%=====================================================================
\section{Ergebnisübersicht}

\begin{center}
\renewcommand{\arraystretch}{1.5}
\begin{tabularx}{\linewidth}{@{}l>{\raggedright\arraybackslash}X>{\raggedright\arraybackslash}p{3.4cm}@{}}
\toprule
\textbf{Nr.} & \textbf{Ergebnis} & \textbf{Werkzeug} \\ \midrule
1a & $T_1 \approx \SI{278.9}{K}$, $v_1 \approx \SI{262}{m/s}$, $\rho_1 \approx \SI{5.0}{kg/m^3}$ & Isentrope, Energiesatz Gas \\
1b & $v \approx \SI{245}{m/s}$, Fehler $\approx \SI{-7}{\percent}$ & Bernoulli, $\rho = \rho_0$ \\
1c & $T_1 \approx \SI{289.1}{K}$, $\rho_1 \approx \SI{4.82}{kg/m^3}$, $\eta_D \approx 0{,}70$ & Energiesatz Gas \\
2 & $\Delta p_{oben} \approx \SI{2206}{Pa}$, $v_{oben}/v_{unten} \approx 1{,}05$ & Hydrostatik \\
7 & $p_2 = \SI{0.5}{bar}$, $p_3 = \SI{1.5}{bar}$, $c_W = \alpha/(1-\alpha)^2 = 2{,}0$ & Bernoulli, Impulssatz \\
8 & $F = \pi R^2\left[(p_1 - p_2) + \rho v_{max}^2/12\right]$ & Impulssatz, $\beta = 4/3$ \\
  & $\dot Q = \pi R^2 \frac{v_{max}}{2}\left[(p_1 - p_2) + \rho v_{max}^2/8\right]$ & Energiesatz, $\alpha = 2$ \\
9 & $h \approx \SI{0.82}{m}$ & Kontinuität, Impulssatz \\
13 & $Y^* \approx \SI{256}{J/kg}$, $P \approx \SI{79}{kW}$, $\dot Q \approx \SI{75}{kW}$, $p_D \approx \SI{4.48}{bar}$, $p_S \approx \SI{1.91}{bar}$ & erweiterte Bernoulli \\
\bottomrule
\end{tabularx}
\end{center}

\end{document}
